\chapter{Plug-ins and extensions}
\label{ch:plugins}

\begin{aside}
A pluggable app lets users add functionality without
rebuilding. \maya{}'s plug-in model is dynamic loading plus a
narrow ABI.
\end{aside}

\section{The ABI}

A plug-in is a shared library that exports a C ABI entry
point:

\begin{lstlisting}
extern "C" {
    void* maya_plugin_init(MayaContext* ctx);
    void  maya_plugin_shutdown(void* state);
}
\end{lstlisting}

\code{init} returns an opaque state pointer; \code{shutdown}
cleans up.

\section{Discovery}

\maya{} scans a plug-in directory at startup:

\begin{lstlisting}
~/.config/myapp/plugins/*.so
\end{lstlisting}

For each, dlopen, find the symbols, call init.

\section{ABI stability}

The ABI surface is tiny and stable across versions. The
plug-in talks to the context via opaque handles; the
context itself can evolve internally.

\section{What plug-ins can do}

\begin{itemize}[leftmargin=*]
  \item Register commands.
  \item Add keybindings.
  \item Subscribe to events.
  \item Provide widgets.
  \item Add menu entries.
\end{itemize}

\section{Sandboxing}

\maya{} does not sandbox plug-ins; they have full process
access. If sandboxing is needed, run plug-ins in a child
process and IPC.

\section{Lua and scripting}

Many TUI apps embed a scripting language (Lua, Python,
JavaScript) for plug-ins. The runtime is heavier but the
barrier to entry is lower.

\maya{} does not ship an embedded interpreter; this is an
opt-in extension.

\section{Plug-in lifecycle}

\code{init} runs at app startup. \code{shutdown} runs at
app exit. Reload at runtime requires unloading the dylib,
which can be brittle.

\section{Recap}

\begin{recap}
\begin{itemize}[leftmargin=*]
  \item Plug-in = shared library with C ABI.
  \item Discovery via scan of a directory.
  \item Narrow ABI for stability.
  \item No sandboxing built-in.
  \item Scripting is opt-in.
\end{itemize}
\end{recap}

\section*{Workshop}

\begin{enumerate}[leftmargin=*]
  \item Build a hello-world plug-in.
  \item Have the plug-in register a custom command.
  \item Try reloading the plug-in at runtime; observe the
    pitfalls.
\end{enumerate}
